\documentclass[11pt,letterpaper]{article}
\usepackage[T1]{fontenc}\usepackage{lmodern}
\usepackage[margin=0.85in]{geometry}\usepackage[hidelinks]{hyperref}
\setlength{\parindent}{0pt}\setlength{\parskip}{6pt}
\title{Learning Note: Editing Record}\author{Learning-Note}\date{2 October 2026}
\begin{document}
\begin{center}
{\Large\bfseries Learning Note: Editing Record}\par
2 October 2026
\end{center}
\section*{Purpose}
Record substantive changes between editing sessions. This is a dated human-readable record, without numbered releases, approval procedures, or duplicate versions. Future edits append a new dated entry and regenerate this PDF.

\section*{2 October 2026 --- Initial note}
\textbf{Created:} \texttt{Electromagnetism-Foundations/main.tex} and its PDF, plus independent learning-plan and editing-record documents in \texttt{Documents/}.

\textbf{Content:} Operational definitions and experimental access to electric and magnetic fields; magnetic moments, electron spin, stationary magnets, and energy; definitions of flux, divergence, and curl; fluid, heat, current, and electromagnetic applications; illustrative calculations; and a limited map of related physical laws.

\textbf{Clarifications from the supplied discussion:}
\begin{itemize}
\item Electric field is external force per signed test charge in a specified frame, with negligible probe disturbance.
\item The magnetic force is perpendicular to velocity and field; the field need not be perpendicular to velocity. One velocity direction cannot determine every field component.
\item Spin is intrinsic angular momentum, not literal mechanical rotation. Atomic beam splitting is distinguished from a free-electron experiment.
\item Fixed-dipole torque, force, and interaction-energy formulas state their approximation and distinguish uniform from nonuniform external fields.
\item Static magnetic fields and intrinsic spin do not require continuous power; resistive coils have losses.
\item Field flux is not necessarily material transport. Zero divergence does not mean zero field. Circular field lines do not necessarily mean nonzero local curl.
\item Definitions and mathematical theorems are distinguished from physical laws. The history is not reduced to a universal definitions-first sequence.
\end{itemize}

\textbf{Build and organization:} Self-contained pdfLaTeX source with embedded references and diagram, suitable for Overleaf. Stable filenames, one topic folder, and only TeX/PDF document files intended for Git. Build artifacts are excluded locally.

\textbf{Validation:} Compiled locally with pdfLaTeX, resolved citations on a second pass, and inspected rendered PDF pages. Recompilation does not require external images, BibTeX, or shell escape.

\textbf{Next:} Begin Maxwell's equations only when requested, following the learning plan. Wave propagation, reflection, and detailed AN720 analysis remain future lessons.
\end{document}
